\documentclass[conference,a4paper,10pt]{IEEEtran}

\usepackage[utf8]{inputenc}
\usepackage[T1]{fontenc}
\usepackage[spanish,es-tabla]{babel}
\usepackage{graphicx}
\usepackage{amsmath}
\usepackage{amssymb}
\usepackage{booktabs}
\usepackage{cite}
\usepackage{microtype}
\usepackage[a4paper, left=0.85cm, right=0.85cm, top=0.75cm, bottom=0.75cm]{geometry}
\usepackage{url}

% ------------------------------------------------
% CONFIGURACIÓN GENERAL
% ------------------------------------------------

\pagestyle{empty}
\pagenumbering{gobble}

\renewcommand{\baselinestretch}{0.93}
\setlength{\parskip}{0pt}
\setlength{\parindent}{1em}
\setlength{\columnsep}{0.16in}

% Figura grande que ocupa las dos columnas
\newcommand{\reportfigwide}[3]{%
\begin{figure*}[!t]
    \centering
    \includegraphics[width=0.98\textwidth]{#1}
    \caption{#2}
    \label{#3}
\end{figure*}
}

% Espaciado de figuras y tablas
\setlength{\textfloatsep}{4pt plus 1pt minus 1pt}
\setlength{\floatsep}{4pt plus 1pt minus 1pt}
\setlength{\intextsep}{4pt plus 1pt minus 1pt}
\setlength{\abovecaptionskip}{2pt}
\setlength{\belowcaptionskip}{-3pt}

% Espaciado de secciones
\makeatletter
\def\section{\@startsection{section}{1}{\z@}
{0.7ex plus 0.2ex minus 0.2ex}
{0.45ex plus 0.1ex}
{\normalfont\normalsize\bfseries}}

\def\subsection{\@startsection{subsection}{2}{\z@}
{0.55ex plus 0.1ex minus 0.1ex}
{0.3ex plus 0.1ex}
{\normalfont\small\bfseries}}
\makeatother

% ------------------------------------------------
% COMANDO PARA FIGURAS
% ------------------------------------------------

\newcommand{\reportfig}[4]{%
\begin{figure}[!htbp]
    \centering
    \includegraphics[width=#1\columnwidth]{#2}
    \vspace{-0.08cm}
    \caption{#3}
    \label{#4}
    \vspace{-0.12cm}
\end{figure}
}

% ------------------------------------------------
% TÍTULO Y AUTORES
% ------------------------------------------------

\title{Trabajo Práctico – Sistema de Gestión de Reservas y Pedidos de un Restaurante}

\author{
\IEEEauthorblockN{Victoria Pérsico, Inés Mestres, Sofía Yabo, Mateo Villegas Leiva}
\IEEEauthorblockA{
Universidad de San Andrés\\
Bases de Datos\\
Primavera / 2026\\
Profesores: Mariano Beiró, Lucas Román
}
}

% =================================================
% DOCUMENTO
% =================================================

\begin{document}


% =================================================
% PARTE V -- CARGA DE DATOS 
% copiar desde aqui
% =================================================

\section{Parte V -- Carga de datos}

Una vez creadas las tablas se poblaron con datos ficticios, Para ello se desarrolló un proceso automático de dos etapas: un programa generador escrito en Python, \texttt{generar\_datos.py}, que produce como salida un script SQL de carga, \texttt{carga\_datos.sql}, el cual se ejecuta luego sobre la base de datos.

\subsection{Volumen de datos generados}

Los datos simulan la actividad del restaurante entre el 1 y el 15 de
septiembre de 2026, e incluyen además reservas futuras hasta el 30 de
septiembre. La Tabla~\ref{tab:volumen} resume la cantidad de tuplas cargadas
en cada relación.

\begin{table}[t]
\centering
\caption{Cantidad de tuplas cargadas por relación.}
\label{tab:volumen}
\footnotesize
\setlength{\tabcolsep}{3.2pt}
\begin{tabular}{lr}
\toprule
Relación & Tuplas \\
\midrule
Aplicacion & 4 \\
CategoriaPlato & 8 \\
Mesa & 20 \\
Plato & 43 \\
Cliente & 30 \\
Reserva & 30 \\
Pedido & 57 \\
\quad PedidoMesa & 17 \\
\quad PedidoTelefonico & 20 \\
\quad PedidoAplicacion & 20 \\
ItemPedido & 186 \\
\bottomrule
\end{tabular}
\end{table}

\subsection{Criterios de consistencia}

La generación no es puramente aleatoria:
\begin{itemize}
    \item \textbf{Especialización total y disjunta.} Cada pedido generado se
    inserta en exactamente una de las tablas \texttt{PedidoMesa},
    \texttt{PedidoTelefonico} o \texttt{PedidoAplicacion}, respetando la
    hipótesis definida en el modelo conceptual.

    \item \textbf{Capacidad de las mesas.} A cada reserva se le asigna la mesa
    de menor capacidad que pueda alojar a la cantidad de personas indicada, de
    modo que nunca se supera la capacidad de la mesa asignada.

    \item \textbf{Disponibilidad de las mesas.} Una misma mesa no se asigna a
    dos reservas dentro del mismo turno, lo que evita superposiciones
    horarias. Se consideraron dos turnos diarios, almuerzo y cena, y se asumió
    que el restaurante no abre los días lunes.

    \item \textbf{Relación entre reservas y pedidos.} Únicamente las reservas
    en estado \texttt{completada}, o las del día en curso que ya comenzaron,
    generan un pedido de mesa, y a lo sumo uno. Las reservas canceladas y las
    futuras no tienen pedido asociado, en línea con la cardinalidad $(0,1)$
    definida en el DEI.

    \item \textbf{Precios históricos.} Se simulo una actualizacion de la
    carta a lo largo del período. El atributo \texttt{precio\_unitario} de cada
    ítem conserva el precio vigente en la fecha del pedido, que por lo tanto
    difiere del \texttt{precio\_actual} del plato.

    \item \textbf{Coherencia temporal de los estados.} Los pedidos de fechas
    pasadas se encuentran en estado \texttt{entregado} o \texttt{cancelado},
    mientras que los del día en curso pueden estar \texttt{pendiente},
    \texttt{en\_preparacion} o \texttt{listo} según el tiempo transcurrido. De
    forma análoga, las reservas pasadas están \texttt{completada} o
    \texttt{cancelada} y las futuras \texttt{pendiente} o \texttt{confirmada}.

\end{itemize}

\begin{table*}[t]
\centering
\caption{Consultas diseñadas, requerimiento funcional atendido y recursos SQL utilizados.}
\label{tab:consultas}
\footnotesize
\setlength{\tabcolsep}{4pt}
\begin{tabular}{clll}
\toprule
\# & Descripción & RF & Recursos SQL \\
\midrule
1 & Reservas realizadas por un determinado cliente, con su mesa y estado & 11 & Junta de tres relaciones \\
2 & Platos más solicitados, en unidades vendidas y en facturación & 12 & Junta, \texttt{GROUP BY}, \texttt{COUNT DISTINCT} \\
3 & Categorías más solicitadas y su participación en la facturación total & 12 & Junta, \texttt{GROUP BY}, función de ventana \\
4 & Pedidos, montos y ticket promedio por mes y canal de origen & 13 & Juntas externas a los subtipos, \texttt{GROUP BY}, \texttt{CASE} \\
5 & Actividad de cada aplicación externa: pedidos, usuarios y facturación & 8, 13 & Junta, \texttt{GROUP BY}, \texttt{COUNT DISTINCT} \\
6 & Clientes con diez o más reservas completadas y su consumo acumulado & 1, 3, 11 & Juntas externas, \texttt{GROUP BY}, \texttt{HAVING} \\
7 & Evolución mensual de las reservas y porcentaje de cancelación & 3 & \texttt{GROUP BY}, \texttt{COUNT} con \texttt{FILTER} \\
8 & Uso de cada mesa y aprovechamiento de su capacidad & 2, 4 & Junta externa, \texttt{GROUP BY}, \texttt{AVG} \\
9 & Variación entre el precio actual y los precios históricos de cada plato & 5, 10 & Junta, \texttt{GROUP BY}, \texttt{HAVING}, \texttt{MIN}/\texttt{MAX} \\
10 & Platos sin ventas recientes o nunca solicitados & 5, 12 & Juntas externas, \texttt{GROUP BY}, \texttt{HAVING} sobre fechas \\
\bottomrule
\end{tabular}
\end{table*}

\subsection{Parámetros de configuración}

El generador expone en su encabezado un conjunto de parámetros que permiten controlar el volumen y la forma de los datos sin modificar su lógica: el período simulado, la cantidad de clientes, las fechas de actualización de precios, la cantidad de reservas y de pedidos. 

\subsection{Ejecución del script de carga}

El script \texttt{carga\_datos.sql} está diseñado para poder ejecutarse reiteradamente. Todas sus sentencias se encuentran dentro de una única transacción, por lo que ante un error la base no queda en un estado intermedio. La primera sentencia es un \texttt{TRUNCATE} con las cláusulas
\texttt{RESTART IDENTITY CASCADE} sobre todas las tablas, que vacía las relaciones y reinicia los contadores asociados a las columnas \texttt{SERIAL}. Las tuplas se insertan mediante sentencias \texttt{INSERT} de múltiples filas, agrupadas en lotes de 500, lo que reduce notablemente el tiempo de carga respecto de una sentencia por tupla.

Dado que los identificadores se insertan de manera explícita, para no alterar el orden cronológico de los pedidos, al final del script se reposicionan las secuencias de las columnas \texttt{SERIAL} mediante la función \texttt{setval}. De esta forma, un \texttt{INSERT} posterior que omita el identificador no genera una colisión con las claves ya cargadas.

La secuencia completa de comandos para construir la base de datos es la siguiente:

\begin{itemize}
    \item \texttt{createdb tp\_restaurante}
    \item \texttt{psql -d tp\_restaurante -f script.sql}
    \item \texttt{psql -d tp\_restaurante -f carga\_datos.sql}
\end{itemize}

Opcionalmente, puede regenerarse el script de carga con otros datos ejecutando
\texttt{python3 generar\_datos.py}.

% =================================================
% PARTE VI -- DISEÑO DE CONSULTAS
% =================================================

\section{Parte VI -- Diseño de consultas}

Sobre la base de datos poblada se diseñaron diez consultas SQL vinculadas con los requerimientos funcionales enunciados en la Parte I. Las consultas se encuentran en los archivos \texttt{consulta<n>.sql}. La Tabla~\ref{tab:consultas} las resume, indicando el requerimiento funcional (RF) que atiende cada una y los recursos del lenguaje SQL utilizados.


\subsection{Criterios de diseño}

Las tres últimas consultas del listado de requerimientos funcionales ---consultar las reservas de un cliente, conocer los platos y categorías más solicitados y obtener información resumida de los pedidos por período y canal de origen--- se corresponden directamente con las consultas 1 a 4. Las restantes amplían esa base hacia aspectos de gestión que el modelo permite responder, como la fidelidad de los clientes, el uso de las mesas o la detección de platos sin salida.

Dos decisiones transversales merecen ser señaladas. En primer lugar, las consultas que calculan montos excluyen los pedidos en estado \texttt{cancelado}, de modo que los importes reflejen el consumo efectivamente realizado. En segundo lugar, el canal de origen de un pedido no es un atributo almacenado, sino que se deduce de la pertenencia del pedido a uno de los subtipos. Para recuperarlo se emplean juntas externas hacia \texttt{PedidoMesa} y \texttt{PedidoTelefonico} combinadas con una expresión \texttt{CASE}, aprovechando que la especialización es total y disjunta: si un pedido no aparece en ninguna de esas dos relaciones, necesariamente proviene de una aplicación.

Asimismo, varias consultas utilizan juntas externas para conservar las tuplas sin correspondencia, cuando esa ausencia es justamente lo que se quiere observar. Es el caso de las mesas que no registran reservas en la consulta 8 y de los platos que nunca fueron solicitados en la consulta 10.


% =================================================
% copiar hasta aqui
% =================================================

\end{document}
